\section{结论}

我们推导了非单态夸克分布（包括非极化、螺旋度与横率分布）的单圈匹配条件。该匹配条件也可以被视为一种因子化：它把准夸克分布与实验中可测量的光锥分布联系起来，从而使人们能够从前者提取后者；前者定义为有限核子动量下不含时间的关联函数，因此可以在格点上计算。为了开展准夸克分布的格点模拟，需要采用横向动量截断作为紫外（UV）发散的调节器。这一紫外调节器的选择会在轴向规范 $A^z=0$ 下产生一个乘以奇异因子的线性发散，
其起源是 Wilson 线的自能。在我们的单圈因子化公式中，该发散可以用主值约定消除。然而，要最终实现所有阶的因子化，重要的是研究这种发散结构在单圈以上如何影响准分布与光锥分布之间的匹配。

\begin{acknowledgments}
我们感谢 J. W. Chen 就 $1/(1-x)^2$ 奇异性进行的有益讨论。本工作部分受到美国能源部基金
DE-FG02-93ER-40762、上海市科学技术委员会基金（No. 11DZ2260700）以及中国国家自然科学基金
（No. 11175114）的资助。
\end{acknowledgments}
